\section{Mid-Term Examination --- Set B}

\LPUPaperHeader{Mid-Term Examination}{B}{90 Minutes}{30 Marks}{CO1 (Units 1), CO2 (Unit 2), CO3 (Unit 3)}

\begin{center}
\sffamily\textbf{\color{AlertRed}Note: Attempt all thirty questions. Each question carries 1 mark. Negative marking: $-0.25$ marks for wrong answers.}
\end{center}

\noindent\textbf{Q1.} \hfill {\small\color{LPUOrange}\textbf{[CO1 | named constant]}}\\[2pt]
By convention, how are named constants usually written in Python?

\begin{enumerate}[label=(\alph*),topsep=2pt,itemsep=1pt,partopsep=0pt,parsep=0pt]
    \item \relax In all lowercase letters
    \item \relax In all uppercase letters
    \item \relax In camelCase
    \item \relax With a leading digit
\end{enumerate}
\vspace{{2mm}}

\noindent\textbf{Q2.} \hfill {\small\color{LPUOrange}\textbf{[CO1 | simultaneous assignment]}}\\[2pt]
What will be the values of \texttt{a} and \texttt{b} after \texttt{a, b = 1, 2}?

\begin{enumerate}[label=(\alph*),topsep=2pt,itemsep=1pt,partopsep=0pt,parsep=0pt]
    \item \relax both a and b are 2
    \item \relax both a and b are 1
    \item \relax a is 2 and b is 1
    \item \relax a is 1 and b is 2
\end{enumerate}
\vspace{{2mm}}

\noindent\textbf{Q3.} \hfill {\small\color{LPUOrange}\textbf{[CO1 | simultaneous assignment]}}\\[2pt]
Which single statement correctly swaps the values of \texttt{x} and \texttt{y} in Python?

\begin{enumerate}[label=(\alph*),topsep=2pt,itemsep=1pt,partopsep=0pt,parsep=0pt]
    \item \relax x == y, y == x
    \item \relax swap(x, y)
    \item \relax x, y = y, x
    \item \relax x = y = x
\end{enumerate}
\vspace{{2mm}}

\noindent\textbf{Q4.} \hfill {\small\color{LPUOrange}\textbf{[CO1 | boolean types]}}\\[2pt]
Which of the following are the two Boolean values in Python?

\begin{enumerate}[label=(\alph*),topsep=2pt,itemsep=1pt,partopsep=0pt,parsep=0pt]
    \item \relax true and false
    \item \relax True and False
    \item \relax 1 and 0 only
    \item \relax Yes and No
\end{enumerate}
\vspace{{2mm}}

\noindent\textbf{Q5.} \hfill {\small\color{LPUOrange}\textbf{[CO1 | numeric data types]}}\\[2pt]
Which numeric data type represents whole numbers in Python?

\begin{enumerate}[label=(\alph*),topsep=2pt,itemsep=1pt,partopsep=0pt,parsep=0pt]
    \item \relax complex
    \item \relax str
    \item \relax float
    \item \relax int
\end{enumerate}
\vspace{{2mm}}

\noindent\textbf{Q6.} \hfill {\small\color{LPUOrange}\textbf{[CO1 | numeric data types]}}\\[2pt]
What is the data type of the value \texttt{3.14} in Python?

\begin{enumerate}[label=(\alph*),topsep=2pt,itemsep=1pt,partopsep=0pt,parsep=0pt]
    \item \relax complex
    \item \relax float
    \item \relax bool
    \item \relax int
\end{enumerate}
\vspace{{2mm}}

\noindent\textbf{Q7.} \hfill {\small\color{LPUOrange}\textbf{[CO1 | operators]}}\\[2pt]
Which operator is used for exponentiation in Python?

\begin{enumerate}[label=(\alph*),topsep=2pt,itemsep=1pt,partopsep=0pt,parsep=0pt]
    \item \relax //
    \item \relax \textasciicircum{}
    \item \relax \%
    \item \relax **
\end{enumerate}
\vspace{{2mm}}

\noindent\textbf{Q8.} \hfill {\small\color{LPUOrange}\textbf{[CO1 | operator precedence and associativity]}}\\[2pt]
What is the result of the expression \texttt{2 + 3 * 4} in Python?

\begin{enumerate}[label=(\alph*),topsep=2pt,itemsep=1pt,partopsep=0pt,parsep=0pt]
    \item \relax 14
    \item \relax 20
    \item \relax 24
    \item \relax 11
\end{enumerate}
\vspace{{2mm}}

\noindent\textbf{Q9.} \hfill {\small\color{LPUOrange}\textbf{[CO1 | augmented assignment operators]}}\\[2pt]
What does the statement \texttt{x += 3} mean in Python?

\begin{enumerate}[label=(\alph*),topsep=2pt,itemsep=1pt,partopsep=0pt,parsep=0pt]
    \item \relax x = 3
    \item \relax x = x + 3
    \item \relax x = x * 3
    \item \relax x = 3 - x
\end{enumerate}
\vspace{{2mm}}

\noindent\textbf{Q10.} \hfill {\small\color{LPUOrange}\textbf{[CO1 | type conversion and rounding]}}\\[2pt]
What does the function \texttt{int(4.9)} return in Python?

\begin{enumerate}[label=(\alph*),topsep=2pt,itemsep=1pt,partopsep=0pt,parsep=0pt]
    \item \relax 4
    \item \relax 5
    \item \relax 4.9
    \item \relax 0
\end{enumerate}
\vspace{{2mm}}

\noindent\textbf{Q11.} \hfill {\small\color{LPUOrange}\textbf{[CO2 | function calls]}}\\[2pt]
How do you call a function named \texttt{greet} in Python?

\begin{enumerate}[label=(\alph*),topsep=2pt,itemsep=1pt,partopsep=0pt,parsep=0pt]
    \item \relax \texttt{greet()}
    \item \relax \texttt{function greet}
    \item \relax \texttt{greet[]}
    \item \relax \texttt{call greet}
\end{enumerate}
\vspace{{2mm}}

\noindent\textbf{Q12.} \hfill {\small\color{LPUOrange}\textbf{[CO2 | type conversion and coercion]}}\\[2pt]
Which function converts a string \texttt{"25"} into an integer?

\begin{enumerate}[label=(\alph*),topsep=2pt,itemsep=1pt,partopsep=0pt,parsep=0pt]
    \item \relax \texttt{int("25")}
    \item \relax \texttt{chr(25)}
    \item \relax \texttt{float("25")}
    \item \relax \texttt{str(25)}
\end{enumerate}
\vspace{{2mm}}

\noindent\textbf{Q13.} \hfill {\small\color{LPUOrange}\textbf{[CO2 | type conversion and coercion]}}\\[2pt]
What is the result of \texttt{float(7)} in Python?

\begin{enumerate}[label=(\alph*),topsep=2pt,itemsep=1pt,partopsep=0pt,parsep=0pt]
    \item \relax \texttt{70}
    \item \relax \texttt{7.0}
    \item \relax \texttt{"7"}
    \item \relax \texttt{7}
\end{enumerate}
\vspace{{2mm}}

\noindent\textbf{Q14.} \hfill {\small\color{LPUOrange}\textbf{[CO2 | math functions]}}\\[2pt]
Which function from the \texttt{math} module returns the square root of a number?

\begin{enumerate}[label=(\alph*),topsep=2pt,itemsep=1pt,partopsep=0pt,parsep=0pt]
    \item \relax \texttt{math.sqrt()}
    \item \relax \texttt{math.pow()}
    \item \relax \texttt{math.root()}
    \item \relax \texttt{math.square()}
\end{enumerate}
\vspace{{2mm}}

\noindent\textbf{Q15.} \hfill {\small\color{LPUOrange}\textbf{[CO2 | math functions]}}\\[2pt]
What does \texttt{math.ceil(4.2)} return?

\begin{enumerate}[label=(\alph*),topsep=2pt,itemsep=1pt,partopsep=0pt,parsep=0pt]
    \item \relax \texttt{4}
    \item \relax \texttt{4.0}
    \item \relax \texttt{4.2}
    \item \relax \texttt{5}
\end{enumerate}
\vspace{{2mm}}

\noindent\textbf{Q16.} \hfill {\small\color{LPUOrange}\textbf{[CO2 | adding new function]}}\\[2pt]
Which keyword is used to define a new function in Python?

\begin{enumerate}[label=(\alph*),topsep=2pt,itemsep=1pt,partopsep=0pt,parsep=0pt]
    \item \relax \texttt{def}
    \item \relax \texttt{func}
    \item \relax \texttt{define}
    \item \relax \texttt{function}
\end{enumerate}
\vspace{{2mm}}

\noindent\textbf{Q17.} \hfill {\small\color{LPUOrange}\textbf{[CO2 | parameters and argument]}}\\[2pt]
In the function definition \texttt{def add(x, y):}, what are \texttt{x} and \texttt{y} called?

\begin{enumerate}[label=(\alph*),topsep=2pt,itemsep=1pt,partopsep=0pt,parsep=0pt]
    \item \relax Parameters
    \item \relax Keywords
    \item \relax Variables
    \item \relax Arguments
\end{enumerate}
\vspace{{2mm}}

\noindent\textbf{Q18.} \hfill {\small\color{LPUOrange}\textbf{[CO2 | recursion and its use]}}\\[2pt]
What is recursion in programming?

\begin{enumerate}[label=(\alph*),topsep=2pt,itemsep=1pt,partopsep=0pt,parsep=0pt]
    \item \relax A function that calls itself
    \item \relax A loop that runs forever
    \item \relax A function with no parameters
    \item \relax A function that returns nothing
\end{enumerate}
\vspace{{2mm}}

\noindent\textbf{Q19.} \hfill {\small\color{LPUOrange}\textbf{[CO2 | recursion and its use]}}\\[2pt]
What is required in every recursive function to prevent infinite recursion?

\begin{enumerate}[label=(\alph*),topsep=2pt,itemsep=1pt,partopsep=0pt,parsep=0pt]
    \item \relax A print statement that displays the current value at each level of the recursive call chain
    \item \relax A base case
    \item \relax A global variable
    \item \relax An infinite loop
\end{enumerate}
\vspace{{2mm}}

\noindent\textbf{Q20.} \hfill {\small\color{LPUOrange}\textbf{[CO2 | conditional expressions]}}\\[2pt]
What is the value of the conditional expression \texttt{"Yes" if 5 > 3 else "No"}?

\begin{enumerate}[label=(\alph*),topsep=2pt,itemsep=1pt,partopsep=0pt,parsep=0pt]
    \item \relax \texttt{"No"}
    \item \relax \texttt{"Yes"}
    \item \relax \texttt{True}
    \item \relax \texttt{False}
\end{enumerate}
\vspace{{2mm}}

\noindent\textbf{Q21.} \hfill {\small\color{LPUOrange}\textbf{[CO3 | length]}}\\[2pt]
What does \texttt{len([10, 20, 30, 40])} return?

\begin{enumerate}[label=(\alph*),topsep=2pt,itemsep=1pt,partopsep=0pt,parsep=0pt]
    \item \relax \texttt{3}
    \item \relax \texttt{4}
    \item \relax \texttt{5}
    \item \relax \texttt{0}
\end{enumerate}
\vspace{{2mm}}

\noindent\textbf{Q22.} \hfill {\small\color{LPUOrange}\textbf{[CO3 | membership]}}\\[2pt]
What does the expression \texttt{3 in [1, 2, 3]} evaluate to?

\begin{enumerate}[label=(\alph*),topsep=2pt,itemsep=1pt,partopsep=0pt,parsep=0pt]
    \item \relax \texttt{False}
    \item \relax \texttt{3}
    \item \relax \texttt{True}
    \item \relax \texttt{None}
\end{enumerate}
\vspace{{2mm}}

\noindent\textbf{Q23.} \hfill {\small\color{LPUOrange}\textbf{[CO3 | operations]}}\\[2pt]
What is the result of \texttt{[1, 2] + [3, 4]}?

\begin{enumerate}[label=(\alph*),topsep=2pt,itemsep=1pt,partopsep=0pt,parsep=0pt]
    \item \relax \texttt{[1, 2, [3, 4]]}
    \item \relax \texttt{[4, 6]}
    \item \relax \texttt{[1, 2, 3, 4]}
    \item \relax \texttt{[3, 4, 1, 2]}
\end{enumerate}
\vspace{{2mm}}

\noindent\textbf{Q24.} \hfill {\small\color{LPUOrange}\textbf{[CO3 | slices]}}\\[2pt]
What is the result of \texttt{[10, 20, 30, 40][1:3]}?

\begin{enumerate}[label=(\alph*),topsep=2pt,itemsep=1pt,partopsep=0pt,parsep=0pt]
    \item \relax \texttt{[20, 30, 40]}
    \item \relax \texttt{[20, 30]}
    \item \relax \texttt{[30, 40]}
    \item \relax \texttt{[10, 20]}
\end{enumerate}
\vspace{{2mm}}

\noindent\textbf{Q25.} \hfill {\small\color{LPUOrange}\textbf{[CO3 | deletion]}}\\[2pt]
Which keyword is used to remove an item from a list by its index?

\begin{enumerate}[label=(\alph*),topsep=2pt,itemsep=1pt,partopsep=0pt,parsep=0pt]
    \item \relax \texttt{del}
    \item \relax \texttt{remove}
    \item \relax \texttt{erase}
    \item \relax \texttt{drop}
\end{enumerate}
\vspace{{2mm}}

\noindent\textbf{Q26.} \hfill {\small\color{LPUOrange}\textbf{[CO3 | list and for loops]}}\\[2pt]
What is printed by \texttt{for x in [1, 2, 3]: print(x)}?

\begin{enumerate}[label=(\alph*),topsep=2pt,itemsep=1pt,partopsep=0pt,parsep=0pt]
    \item \relax The list \texttt{[1, 2, 3]}
    \item \relax The sum \texttt{6}
    \item \relax The length \texttt{3}
    \item \relax Each element \texttt{1}, \texttt{2}, \texttt{3} on separate lines
\end{enumerate}
\vspace{{2mm}}

\noindent\textbf{Q27.} \hfill {\small\color{LPUOrange}\textbf{[CO3 | list parameters and nested list]}}\\[2pt]
What is a nested list?

\begin{enumerate}[label=(\alph*),topsep=2pt,itemsep=1pt,partopsep=0pt,parsep=0pt]
    \item \relax A list with only numbers
    \item \relax A list that contains one or more lists as elements
    \item \relax A list stored in a file
    \item \relax A list that cannot be changed
\end{enumerate}
\vspace{{2mm}}

\noindent\textbf{Q28.} \hfill {\small\color{LPUOrange}\textbf{[CO3 | tuples]}}\\[2pt]
Which symbol pair is typically used to create a tuple?

\begin{enumerate}[label=(\alph*),topsep=2pt,itemsep=1pt,partopsep=0pt,parsep=0pt]
    \item \relax Parentheses \texttt{()}
    \item \relax Square brackets \texttt{[]}
    \item \relax Curly braces \texttt{\{\}}
    \item \relax Angle brackets \texttt{<>}
\end{enumerate}
\vspace{{2mm}}

\noindent\textbf{Q29.} \hfill {\small\color{LPUOrange}\textbf{[CO3 | operations on tuples]}}\\[2pt]
What key property makes tuples different from lists?

\begin{enumerate}[label=(\alph*),topsep=2pt,itemsep=1pt,partopsep=0pt,parsep=0pt]
    \item \relax Tuples are immutable and cannot be changed after creation
    \item \relax Tuples use square brackets
    \item \relax Tuples can only hold strings
    \item \relax Tuples cannot be indexed
\end{enumerate}
\vspace{{2mm}}

\noindent\textbf{Q30.} \hfill {\small\color{LPUOrange}\textbf{[CO3 | creating dictionary]}}\\[2pt]
Which of the following correctly creates a dictionary?

\begin{enumerate}[label=(\alph*),topsep=2pt,itemsep=1pt,partopsep=0pt,parsep=0pt]
    \item \relax \texttt{\{"a", 1, "b", 2\}}
    \item \relax \texttt{["a": 1, "b": 2]}
    \item \relax \texttt{\{"a": 1, "b": 2\}}
    \item \relax \texttt{("a", 1, "b", 2)}
\end{enumerate}
\vspace{{2mm}}

\newpage
\subsection*{Official Answer Key \& Analytical Explanations --- Set B}

\begin{center}
\begin{tabularx}{0.96\textwidth}{c c X c c X}
\toprule
\textbf{Q\#} & \textbf{Ans} & \textbf{Core Concept / Rule} & \textbf{Q\#} & \textbf{Ans} & \textbf{Core Concept / Rule} \\
\midrule
1 & \textbf{D} & Communicates instructions to CPU & 16 & \textbf{A} & Conditional if keyword begins block \\
2 & \textbf{B} & High-level interpreted language & 17 & \textbf{B} & Modulo parity condition \\
3 & \textbf{C} & Enforces indentation block syntax & 18 & \textbf{C} & Ternary expression evaluation \\
4 & \textbf{B} & Dynamic type inference & 19 & \textbf{D} & Nested loop iteration product \\
5 & \textbf{A} & Guido van Rossum in 1991 & 20 & \textbf{A} & Break terminates loop immediately \\
6 & \textbf{C} & Case-sensitive naming rules & 21 & \textbf{B} & Continue advances to next step \\
7 & \textbf{B} & Reserved keyword cannot be identifier & 22 & \textbf{D} & Loop else runs without break \\
8 & \textbf{C} & Variable reassignment to new address & 23 & \textbf{C} & Mersenne Twister seed state \\
9 & \textbf{A} & Binds object 5 to identifier x & 24 & \textbf{A} & Positional argument binding \\
10 & \textbf{D} & IEEE 754 float conversion & 25 & \textbf{B} & Avoid mutable default argument \\
11 & \textbf{B} & Floor division round toward -inf & 26 & \textbf{D} & Base case stops infinite recursion \\
12 & \textbf{A} & Simultaneous tuple unpacking & 27 & \textbf{C} & Immutable sequence slice \\
13 & \textbf{D} & Tuple swap without temp variable & 28 & \textbf{B} & Negative step reverses string \\
14 & \textbf{B} & Bitwise left-shift identity & 29 & \textbf{A} & List append amortized O(1) \\
15 & \textbf{C} & Operator precedence ** over * & 30 & \textbf{D} & Dictionary O(1) hash table lookup \\
\bottomrule
\end{tabularx}
\end{center}
\vspace{{3mm}}
